\documentclass[11pt,a4paper]{article}
\usepackage[margin=24mm]{geometry}
\usepackage{amsmath,amssymb,array}
\input{glyphtounicode}
\pdfgentounicode=1
\usepackage[hidelinks]{hyperref}
\hypersetup{hypertexnames=false}
\setlength{\parindent}{0pt}
\setlength{\parskip}{5pt}
\newcommand{\pos}[1]{[#1]_+}
\begin{document}
\begin{center}
{\Large M2 joined opening: admission before recovery}\\[4pt]
15 October 2026 --- analytical audit and a terminal-priority gate
\end{center}

\textbf{Status.} The user approved the three-part joined review and conditional
E slowing. M2 remains selected. This note checks a shared service/allocation
interface and stops at M2-A1, a possible narrow terminal-priority amendment.
Neither alternative nor a complete M2 constitutive law is approved. No
traffic integration, numerical gradient, calibration or optimisation occurs.
All earlier mathematical sources/PDFs remain unchanged.

\section*{1. Retained physical inputs and the literal terminal rule}

Retain H1's entrance histories $F_j$, minimum-travel eligibility $P_j$,
finite terminal length $b$, jam density $k_j$, common advice $v$ and known
empty initial state. Write
\begin{equation}
 P_j(t)=F_j(t-b/v),\qquad q_j=P_j-Y_j,\qquad N_j=F_j-Y_j,
 \quad j=E,A,B. \tag{1}
\end{equation}
Cumulative mass/stock is in vehicle-equivalents; departure, demand and
service rates below are in vehicle-equivalents per second. Retain
\begin{equation}
 C(v)=\frac{wk_jv}{w+v},\qquad 0\leq\mu(t)\leq R\leq C(v_-),
 \qquad v\in[v_-,v_+]. \tag{2}
\end{equation}
$\mu$ is the one effective shared ceiling. Its M2 update is not specified
by (2). Eligible lane demand has the existing interpretation
\begin{equation}
 d_j=\begin{cases} C(v),&q_j>0,\\ \dot P_j,&q_j=0.\end{cases}
 \tag{3}
\end{equation}
The inherited strict allocator is
\begin{equation}
 g_A=\min\{d_A,\mu\},\qquad
 g_E=\min\{d_E,\pos{\mu-g_A}\},\qquad \dot Y_j=g_j. \tag{4}
\end{equation}

\textbf{Strict-priority consequence.} Whenever $q_A>0$, (2)--(3) give
$d_A=C(v)\geq\mu$, hence
\begin{equation}
 q_A>0\quad\Longrightarrow\quad g_A=\mu,\quad g_E=0. \tag{5}
\end{equation}
The consequence holds for every admissible service history and every
recovery memory. Lowering $\mu$ cannot by itself create an E allocation.
This is temporary exclusion during A backlog, not proof of permanent
starvation of a finite cohort. It does not rule out E entry while A is
unqueued or A delay caused by earlier E entry.

\textbf{Interpretation limit.} Equations (3)--(5) average within-lane
service; they do not identify physical gaps. Giving E passage during
$q_A>0$ by a new demand adapter would change this contract just as changing
(4) would. It is not a local transcription correction.

\newpage
\section*{2. A concrete alternative allocation, not a recovery law}

The merge template in
\href{https://arxiv.org/pdf/1310.2660}{Jin, Gan and Lebacque, section 6,
equations (15)} separates downstream service and priority. Reuse only its
allocation shape, with a proposed E priority weight $0\leq\beta<1$:
\begin{align}
 g_E&=\min\{d_E,\max\{\mu-d_A,\beta\mu\}\},\nonumber\\
 g_A&=\min\{d_A,\max\{\mu-d_E,(1-\beta)\mu\}\}. \tag{6}
\end{align}
This is a proposed M2 adaptation of that allocation; the source's
phenomenological capacity-drop law is not imported. A weight favoring A
would use $\beta<1/2$, but no value or calibration is selected. The weight
is a shared physical/model parameter, not another control or an observer
privilege. At $\beta=0$, (6) reproduces (4).

For all nonnegative demands and service,
\begin{equation}
 0\leq g_j\leq d_j,\qquad
 g_A+g_E=\min\{d_A+d_E,\mu\}. \tag{7}
\end{equation}
To verify, first take $d_A+d_E\leq\mu$: both minima equal their own
demands. Otherwise, if $d_E\leq\beta\mu$, E gets its demand and A the
remainder; if $d_A\leq(1-\beta)\mu$, reverse those roles. In the remaining
case each receives its weighted share. These cases exhaust the domain,
including equalities and $\mu=0$.

If both queues are positive, $d_A=d_E=C(v)\geq\mu$, so
\begin{equation}
 (g_A,g_E)=((1-\beta)\mu,\beta\mu). \tag{8}
\end{equation}
If only A is queued, $g_E=\min\{d_E,\beta\mu\}$; an E stream with positive
demand can therefore enter before the A queue empties. Unused E allocation
returns to A. A lone stream can use the whole ceiling. Thus $\beta$ does
not impose a fixed completed merging proportion outside contention.

\textbf{Hand case, not a rollout.} At $\mu=600$ vehicle-equivalents/hour
and saturated eligible A/E demands, strict priority gives $(600,0)$;
illustrative $\beta=1/4$ gives $(450,150)$. These are exact allocations,
not proposed traffic parameters or empirical throughput. The common total
is unchanged. If total demand is 400, (7) passes 400 under either rule;
no artificial queue or slowing is implied by (6).

\textbf{Scope.} Changing (4) to (6) with $\beta>0$ changes terminal
right-of-way/service allocation. It leaves H1's own-lane-first lateral
rule, cutoff, storage, advice and local completion untouched. The allocation
alone specifies neither an acceptable physical gap, crossing speed,
maneuver friction nor recovery. It is not a complete M2 law.

\newpage
\section*{3. Kinematic and recovery guardrails}

For a triangular fundamental diagram, a congested branch compatible with
a rate $0<f<C(v)$ is
\begin{equation}
 k_{\rm c}(f)=k_j-f/w,\qquad
 u_{\rm c}(f)=\frac{wf}{wk_j-f},\qquad C(u_{\rm c}(f))=f. \tag{9}
\end{equation}
The free branch at the same rate has speed $v$. A queue/regime can motivate
examining (9); the lateral insertion trajectory does not follow uniquely
from a point count. No actual speed is assigned to absent traffic.

The source
\href{https://repository.tudelft.nl/file/File_c424c1b3-7374-436b-a6a7-abdf69063aa6}{Leclercq et al., section 3.2}
uses the congested-entry regime for its flow/speed relation. It does not
prove a universal empty-start point-entry map. Its mean-capacity process
and our finite transient account remain distinct.

\textbf{Why an automatic A feedback needs care.} Suppose a tentative rule
uses (9) for queued A and takes the fresh isolated loss amplitude to be
$D_A(0)=R-C(u_A)$. Under strict priority, (5) gives
\begin{equation}
 q_A>0,\quad u_A=u_{\rm c}(g_A),\quad g_A=\mu
 \quad\Longrightarrow\quad D_A(0)=R-\mu. \tag{10}
\end{equation}
The supposed new A-loss amplitude is exactly the shared service reduction
already acting. Adding a fresh independent recovery copy for every such
follower needs another physical justification; observing low A speed alone
does not supply it. Under (8) it would even include the part of service
allocated to E. This diagnostic does not rule out a separately justified
A startup/acceleration mechanism or prove every possible A model duplicates
loss. No such extra source is selected.

\textbf{Shared-account contract.} Let $Z$ denote sufficient predicted
aggregate state/history, not traffic mass, and let a future qualified law
produce $\mu=\mathcal M(Z,\text{current state})$. Admissions can update
$Z$ using their amount/type and justified kinematics. A single output
$\mu$ is compatible with richer internal memory. Neither an unproved
one-state reduction nor independent constraints per past entry are assumed.
An E passage with no actual disturbance must not acquire a fixed slow-entry
penalty merely because it is E.

\textbf{Still open.} Conditional admission/opportunity, actual departure
speed, recovery overlap and causal closure remain unformed. Neither (6)
nor a larger memory automatically supplies these. Initial ordinary
underload may still be a benign regime; a compulsory interaction must be
specified before claiming the model generates a consequential episode.
No exclusive slowing-origin choice is reinstated.

\newpage
\section*{4. Joined interfaces that can already be checked conditionally}

Assume departures are conservative, absolutely continuous, minimum-travel
eligible and demand/service feasible. The fixed finite-terminal domain
\begin{equation}
 H=k_jb>M,\qquad 0\leq\dot F_j\leq C(v),\qquad F_j\leq M \tag{11}
\end{equation}
protects H1's handoff under either allocation. Before $b/w$, vacancy is
at least $(k_jw-C(v))t>0$; afterwards,
\begin{equation}
 V_j(t)=Y_j(t-b/w)+H-F_j(t)\geq H-M>0. \tag{12}
\end{equation}
Physical stock remains $N_j=F_j-Y_j$; $Z$ is never added to stock. These
are sufficient interface bounds, not a general upstream spillback solver.

For the uniform initial footprint $[0,\ell]$ on the road of length $L$,
complete local cohort time has the unchanged conservative identity
\begin{equation}
 J(v)=\frac{M(L-\ell/2)}v+
       \int_0^\infty(q_A+q_E+q_B)\,dt. \tag{13}
\end{equation}
It follows by writing $M-\sum_jY_j=(M-\sum_jP_j)+\sum_jq_j$ and
integrating each travel-eligibility delay. A and E queue areas both count.
Post-crossing recovery may affect following departures but contributes no
second post-crossing vehicle-time to the local objective.

\textbf{A sufficient clearance criterion, not a new M2 proof.} Set
$Q=q_A+q_E$ and $t_*=L/v$, after which all traffic is eligible and no new
external traffic arrives. B retains its independent admitted domain and
$Y_B=P_B$, so B is clear by $t_*$. If a candidate guarantees a finite total loss
budget
\begin{equation}
 0\leq\mu\leq R,\qquad
 \int_{t_*}^{\infty}(R-\mu(t))\,dt\leq\Lambda<\infty, \tag{14}
\end{equation}
then work conservation (7), or (4), gives $\dot Q=-\mu$ whenever $Q>0$.
If it had not cleared by $T>t_*$,
\begin{equation}
 Q(T)=Q(t_*)-R(T-t_*)+
       \int_{t_*}^T(R-\mu(t))dt
 \leq Q(t_*)-R(T-t_*)+\Lambda. \tag{15}
\end{equation}
Thus clearance occurs by $t_*+[Q(t_*)+\Lambda]/R$. The argument permits
temporary zero service and assumes no artificial positive floor. M2 must
establish (14) or another valid clearance argument for its exact law;
M1's finite-kernel proof does not transfer automatically.

\textbf{Limits and gate.} Neither a complete coupled-law uniqueness proof
nor a reachable consequential MP6 domain is supplied. Known initial state
plus a uniquely autonomous law would permit the familiar ending-statistic
initialization, but that conditional statement does not establish decision
relevance or a useful gradient. M2-A1 asks whether to retain (4) or reopen
terminal allocation for positive-weight (6). Both preserve conditional E
slowing and the model--SUMO S-P--SUMO S-L ladder. Freeze no new priority,
demand domain, information channel or recovery law at this audit.
\end{document}
